\section{Conclusion}
\label{sec:conclusion}

We set out to test whether a 3-round adversarial retraining loop, with
or without an LLM rewriting agent, can harden a BiomedBERT detector
against AI-generated biomedical abstracts. On a controlled
\(3{,}000\)-sample corpus, evaluated under fixed seeds and
hyperparameters on a single Apple M3~Max laptop in
\(10\,\mathrm{h}\,01\,\mathrm{m}\) of wall-clock, the answer is no — and
the negative result is informative.

Three findings stand out. First, the choice of \emph{generator family}
explains roughly \(140\times\) more variance in evasion rate than the
choice of retraining schedule or rewriting agent: BioMistral fakes evade
at \(0.7\%\), SeqGAN fakes at \(98.7\%\), RAID's mixed LLMs at
\(98.3\%\), all on the same detector. Second, three rounds of
adversarial retraining produced no measurable improvement on either
in-family (Condition~C, D) or out-of-family (Condition~B) fakes; the
common best-validation-AUC checkpoint selector silently rolled
Conditions~C and~D back to the baseline weights, which is the source of
their byte-identical metrics. Third, a previous \(500\)-sample evasion
curve of \(0.74 \to 0.54\) in Condition~B did not reproduce at
\(3{,}000\) scale and is best explained as small-sample noise within the
\(\pm 0.062\) Wilson CI on \(n = 250\).

The actionable takeaways for downstream practitioners are concrete: any
deployed AI-text detector should publish a
generator-family evasion matrix as part of its documentation; the
default best-validation-AUC checkpoint selector should be replaced by an
adversarial-aware selector when the baseline is at ceiling; and
sub-thousand-sample evasion claims should be treated as exploratory
until they are replicated at a scale where Wilson CIs do not absorb the
effect being claimed. We position BioShield as a biomedical specialist
detector with documented cross-domain failure, release the full code,
data, metrics, and one-command rerun script at
\url{https://github.com/Rayudualapati-25/BioShield}, and treat the
multi-seed replication, LoRA-hardened generator, and family-cross-check
matrix as the highest-priority follow-ups for the paper revision.
